\begin{topic}[id=indoor_context,
             difficulty={Mittel},
             type={Systemanalyse + Vergleich},
             prerequisites={Grundlagen Sensorik; Interesse an Lokalisierung und kontextsensitiven Systemen},
             practice={optional},
             owner={Dozententeam},
             updated={2026-04-09},
             tags={ Indoor-Lokalisation, Kontext, Sensorfusion, Privacy, Edge }]{Indoor-Lokalisation \& kontextsensitive Systeme}

\TopicDescription{Kontextsensitive Systeme kombinieren Lokalisierung, Sensordaten und Regeln, um adaptive Dienste zu liefern (z.~B. Gebäude, Logistik). Wir diskutieren Verfahren, Datenschutz by Design und robuste Edge-Architekturen. Ziel: Konzepte sicher planen und Grenzen transparent machen.}

\TopicLeitfragen{\item Welches Verfahren passt zu welchem Genauigkeits-/Kosten-Ziel?
\item Wie binde ich Privacy by Design ein?
\item Wie bleibt das System wartbar?
}
\TopicScope{\item Kein Großprojekt-Framework.
\item Keine umfangreiche App-Entwicklung.
}
\TopicPractice{Optionale Praxisidee: Architektur-Skizze mit Datenschutz-Konzept und Ausfallszenarien.}

\TopicKeywords{Indoor-Lokalisation, Kontext, Sensorfusion, Privacy, Edge}

\nocite{contextAwareSurvey,privacyByDesign}
\end{topic}
